% R15 component; only input paths and enumerated location/grammar prose are edited.
% Source: revisions/2026-10-04-r15/manuscript/new_theory.tex; commit 1cb3cc9efe135966ec228dfaf51c6841a6dfc97c.
\subsection{An operational evaluation--improvement account}
\label{sec:r15mechanism}
The costate in a rollout is the derivative of a specified continuation
problem. Its approximation error has to be measured where the new policy
uses that continuation. We now make both objects explicit. The construction
uses the capital economy, including its dense production interaction,
common and idiosyncratic shocks, consumption curvature, and terminal
dispersion preference. It also accommodates a critic fitted on a coarser
grid than the one used for assessment.

Fix the assessment grid $t_k=kh$, $Nh=T$. Write $\pi_k=\bar M_k/h$,
where $\bar M_k$ is the disclosed numerical representative of
$\int_{t_k}^{t_{k+1}}m_0(t)dt$. In the following finite problem, these
representatives are fixed real numbers. They specify a feasible,
state-independent reference. The internal transition and transformed
reward are
\begin{align}
 Z_{k+1}&=Z_k+h\{c\bm1+f(Z_k)-a_k\}
                  +\Sigma\sqrt h\,\operatorname{clip}_{z}(G_k),
       \qquad G_k\sim N(0,I_{d+1}),\label{eq:r15finite}\\
 \ell_k(z,a)&=B_k\overline{f(z)}+A_k\overline{\log a}
       -B_k\bar a-\frac{\zeta A_k}{2}\bar a^2,
       \qquad g_N(z)=-\chi e^{-\rho T}\|\Pi z\|_d^2.
       \label{eq:r15reward}
\end{align}
Here $A_k,B_k$ are the integrated weights defined above; all discounting
is included in them. Let $\mathcal R_h(a)$ be the expected sum of these
rewards and $g_N$. The candidate may depend on its entire observed
innovation history. Its actions and the reference belong to the same
declared convex action box $\mathcal A_k$ at each date.

Let $V_k^\pi$ be the value of this finite problem under $\pi$, and define
its continuation after choosing the deterministic drift by
\begin{equation}
 S_k^\pi(x)=\E\big[V_{k+1}^\pi
       \{x+\Sigma\sqrt h\operatorname{clip}_{z}(G_k)\}\big],
 \qquad q_k^\pi=dD_xS_k^\pi.
 \label{eq:r15postdecision}
\end{equation}
The argument $x$ is the conditional mean before the current innovation.
The current reward $\ell_k$ is excluded from $S_k^\pi$; subsequent
rewards start at $k+1$ on the same global grid. In particular,
$q_{N-1}^\pi(x)=-2\chi e^{-\rho T}\Pi x$. Smoothness of $f$, bounded
actions, finite moments, and the quadratic terminal reward justify the
derivatives and expectations in this finite problem.

After fitting, freeze an arbitrary continuously differentiable scalar
predictor $\widehat S_k$ and its costate $\widehat q_k=dD\widehat S_k$.
Set
\[
 x_b(z)=z+h\{c\bm1+f(z)-b\},\qquad
 \widehat Q_k(z,b)=\ell_k(z,b)+\widehat S_k\{x_b(z)\}.
\]
For $U$ uniform on $[0,1]$, independent of the candidate's past, the
\emph{action bridge} is
$X_{k,U}=x_{\pi_k}(Z_k)+U\{x_{a_k}(Z_k)-x_{\pi_k}(Z_k)\}$.
Define the following quantities under the candidate's internal-state law:
\begin{align}
 \mathcal M_h&=\sum_{k=0}^{N-1}\E\big[
       \widehat Q_k(Z_k,a_k)-\widehat Q_k(Z_k,\pi_k)\big],\notag\\
 \mathcal A_h&=\sum_{k=0}^{N-1}h\E\|a_k-\pi_k\bm1\|_d^2,
 &\mathcal E_h&=\sum_{k=0}^{N-1}h\E
       \|\widehat q_k(X_{k,U})-q_k^\pi(X_{k,U})\|_d^2.
 \label{eq:r15accounts}
\end{align}
$\mathcal M_h$ records the improvement predicted by the fitted scalar
continuation, $\mathcal A_h$ records the size of the implemented change,
and $\mathcal E_h$ measures the error that can affect this particular
change. Training-state accuracy alone does not identify $\mathcal E_h$.

\begin{proposition}[Finite Bellman bridge]
\label{prop:r15bellmanbridge}
If the displayed quantities are integrable, then
\begin{equation}
 \mathcal R_h(a)-\mathcal R_h(\pi)
       \geq\mathcal M_h-\sqrt{\mathcal A_h\mathcal E_h}.
 \label{eq:r15bridge}
\end{equation}
The assertion holds for every feasible candidate. It requires neither
exact maximization of $\widehat Q$ nor a continuous-time costate
approximation. A predictor fitted on another mesh enters
$\mathcal E_h$ with its full discrepancy from the assessment-grid
continuation.
\end{proposition}

The proposition follows from the finite-horizon performance-difference
identity and a line integral of the costate error along each action
bridge. It is useful even when action optimization is incomplete. The
next result additionally quantifies that incompleteness with a finite,
checkable calculation.

\begin{proposition}[A feasible-actor account]
\label{prop:r15actorgap}
Suppose $\widehat S_k$ is twice continuously differentiable and
$dD^2\widehat S_k(x)\preceq\Lambda_k I_d$ everywhere. Let
$m_{+,k}$ be the upper endpoint of $\mathcal A_k$, and put
\[
 \Gamma_k=h^2\Lambda_k-A_k/m_{+,k}^2,\qquad
 g_{k,i}=A_k\{a_{k,i}^{-1}-\zeta\bar a_k\}
                    -B_k-h\widehat q_{k,i}\{x_{a_k}(Z_k)\}.
\]
An upper bound on the actor's $\widehat Q_k$ maximization gap is
\begin{equation}
 \delta_k=\max_{b\in\mathcal A_k}
     \left\{\langle g_k,b-a_k\rangle_d
                      +\frac{\Gamma_k}{2}\|b-a_k\|_d^2\right\}.
 \label{eq:r15actorgap}
\end{equation}
This is a separable scalar calculation. If $\Gamma_k<0$, its maximizer
is the coordinate projection of $a_k-g_k/\Gamma_k$. If
$\Gamma_k\geq0$, one compares the two endpoints in each coordinate.
Thus \eqref{eq:r15actorgap} remains valid when global strong concavity
of the learned continuation problem cannot be verified.

If $\mu_0=\min_k(-\Gamma_k/h)>0$ is verified and
$\mathcal D_h=\sum_k\E\delta_k$, then
\begin{equation}
 \mathcal R_h(a)-\mathcal R_h(\pi)
       \geq\tfrac12\mathcal M_h-\mathcal D_h
                            -2\mathcal E_h/\mu_0.
 \label{eq:r15quadratic}
\end{equation}
For exact $\widehat Q$ maximizers, $\mathcal D_h=0$ and the coefficient
of $\mathcal E_h/\mu_0$ improves from two to one. These conclusions
apply to the actual feasible action; a small training loss is not
substituted for $\delta_k$.
\end{proposition}

\paragraph*{Independent occupation and continuation samples.}
Draw an independent candidate trajectory, $K$ uniformly from
$\{0,\ldots,N-1\}$, and $U$ uniformly from $[0,1]$. At its bridge
state, draw two independent continuation banks $Z_1,Z_2$ with
conditional mean $q_K^\pi(X_{K,U})$. A bank may contain antithetic
pairs; independence is required between the two banks, not between
the members of an antithetic pair. The statistics
\begin{align}
 M&=N\{\widehat Q_K(Z_K,a_K)-\widehat Q_K(Z_K,\pi_K)\},
 &A&=T\|a_K-\pi_K\bm1\|_d^2,\notag\\
 E&=T\langle\widehat q_K(X_{K,U})-Z_1,
                    \widehat q_K(X_{K,U})-Z_2\rangle_d,
 &D&=N\delta_K
 \label{eq:r15statistics}
\end{align}
have expectations $\mathcal M_h,\mathcal A_h,\mathcal E_h,
\mathcal D_h$. Individual values of $E$ may be negative. Replacing
them by their positive parts would change the estimand. Moreover,
$E^{\rm raw}=T\|Z_1-Z_2\|_d^2/2$ estimates the integrated mean-square
error of one raw bank on this same occupation law. A confidence bound
for $E-E^{\rm raw}$ therefore assesses the critic's reusable costate
prediction against a direct bank with the declared number of rollouts.
It concerns this evaluation subproblem; the direct method comparison
continues to use the returned policies' independently assessed payoffs.

The supplement gives explicit clipping ranges and expectation
allowances for these statistics. They exploit the common quadratic
terminal component and antithetic cancellation. They do not use an
empirical maximum as a population bound. Independent, equally sized
samples for every member of a declared finite stream distribution may
be pooled using the maximum proved range and the mean numerical and
tail allowances. The resulting target is the uniform mean over all
declared streams.

\begin{theorem}[From fitted continuation to continuous economic gain]
\label{thm:r15mechanism}
Condition on the fitted objects and their selection. Suppose a fresh,
simultaneous assessment gives $L_M\leq\mathcal M_h$,
$\mathcal A_h\leq U_A$, $\mathcal E_h\leq U_E$, and, when used,
$\mathcal D_h\leq U_D$. Include the disclosed arithmetic and
state-representation errors in these endpoints. Define the deterministic
holding deficit
\[
 D_{\rm hold}=\sum_k\left\{C_k-A_k\log\pi_k+B_k\pi_k
                                      +\zeta A_k\pi_k^2/2\right\}\geq0.
\]
For the sampled continuous-economy implementation of
Proposition~\ref{prop:r11transfer}, let $E_{a,h}^{\rm all}$ be its
saved paired payoff-transfer allowance. Then, on that same event,
\begin{equation}
 J(a)-J(m_0)\geq
 L_M-\sqrt{\max(0,U_A)\max(0,U_E)}
                  -D_{\rm hold}-E_{a,h}^{\rm all}.
 \label{eq:r15economicgain}
\end{equation}
When the positive curvature modulus in
Proposition~\ref{prop:r15actorgap} is verified, one may also use
$L_M/2-U_D-2\max(0,U_E)/\mu_0$ before subtracting the same two
allowances. Taking the larger of these two bounds spends no additional
probability. Subtracting a valid lower gain bound from the analytical
reference regret bound gives a bound against all adapted feasible
controls in the original economy.
\end{theorem}

This account separates three observable economic requirements: a
continuation that predicts a useful feasible change, an error small
enough under the candidate's occupation law, and a sufficiently accurate
implementation. The finite Bellman identity is exact; the transfer of
payoffs to the diffusion is quantitative. Neither equality of discrete
and continuous costates nor global concavity of the full adapted-control
problem is assumed.
